\chapter{Wiederholung aus der Funktionalanalysis}

\begin{definition}
    Sei $\Aa$ eine nichtleere Menge und $\cdot : \Aa \times \Aa \to \Aa$ eine Abbildung. Das Tupel $(\Aa, \cdot)$ heißt \emph{Algebra}, wenn
    \begin{itemize}
        \item $\Aa \neq \{ 0 \}$,
        \item $\Aa$ ein $\C$-Vektorraum, und
        \item $\cdot$ assoziativ und bilinear ist.
    \end{itemize}
    Wir schreiben auch nur $\Aa$ wenn $\cdot$ aus dem Kontext klar ist.
\end{definition}

\begin{example}
    Für ein Maß $\mu$ ist $\ell^\infty(\mu, \C)$ mit punktweiser Multiplikation eine Algebra.
\end{example}

\begin{definition}
    Sei $\Aa$ eine Algebra versehen mit einer Norm $\Vert \cdot \Vert$. Erfüllt diese $\Vert a \cdot b \Vert \leq \Vert a \Vert \Vert b \Vert$ für alle $a, b \in \Aa$, so heißt $(\Aa, \Vert \cdot \Vert)$ \emph{normierte Algebra}. Wir schreiben auch nur $\Aa$, wenn die Norm aus dem Kontext klar ist.

    Ist eine normierte Algebra $\Aa$ vollständig als Banachraum, so heißt $\Aa$ \emph{Banachalgebra}.
\end{definition}

\begin{example}
    Sei $X$ ein Banachraum und bezeichne $L_b(X)$ die Menge aller beschränkten linearen Abbildungen. Dann ist $L_b(X)$ versehen mit Komposition und der Operatornorm eine Banachalgebra.
\end{example}

\begin{definition}
    Sei $(\Aa, \cdot)$ eine Algebra, und ${}^* : \Aa \to \Aa$ involutorisch, konjugiert linear mit $(a \cdot b)^* = b^* \cdot a^*$. Dann heißt $(\Aa, {}^*)$ \emph{$\ast$-Algebra}. Wir schreiben auch nur $\Aa$, wenn $\cdot$ und ${}^*$ aus dem Kontext klar sind.
\end{definition}

\begin{example}
    Der Raum der komplexen Polynome $\C[z]$ ist versehen mit der Abbildung ${}^* : p(z) \mapsto \overline{p(z)}$ eine $\ast$-Algebra.
\end{example}

\begin{definition}
    Sei $(\Aa, \cdot, {}^*)$ eine $\ast$-Algebra und $\Aa$ eine normierte Algebra. Ist ${}^*$ isometrisch, so heißt $\Aa$ \emph{normierte $*$-Algebra}. Entsprechend definiert man \emph{Banach-$*$-Algebra}, wir schreiben auch $B^*$-Algebra.
\end{definition}

\begin{definition}
    Eine \emph{$C^*$-Algebra} ist eine $B^*$-Algebra $\Aa$, in der $\Vert a^* a \Vert = \Vert a \Vert^2$ für alle $a \in \Aa$ gilt.
\end{definition}

\begin{lemma}
    Sei $\Aa$ eine Banachalgebra, ${}^* : \Aa \to \Aa$ konjugiert linear und involutorisch und für alle $a, b \in \Aa$ gelte $(a \cdot b)^* = b^* \cdot a^*$. Gilt für alle $a \in \Aa$, dass $\Vert a \Vert^2 \leq \Vert a a^* \Vert$, so ist $(\Aa, {}^*)$ sogar eine $C^*$-Algebra.
\end{lemma}

\begin{proof}
    Sei $a \in \Aa, a \neq 0$ beliebig, dann gilt
    $$ \Vert a \Vert^2 \leq \Vert a a^* \Vert \leq \Vert a \Vert \Vert a^* \Vert. $$
    Wegen $\Vert a \Vert \neq 0$ folgt $ \Vert a \Vert \leq \Vert a^* \Vert $ für alle $a \in \Aa$. Mit $a' \coloneqq a^*$ folgt mit derselben Argumentation Gleichheit. Weiters folgt
    \begin{equation*}
        \Vert a \Vert^2 \leq \Vert a a^* \Vert \leq \Vert a \Vert \Vert a^* \Vert = \Vert a \Vert^2. \qedhere
    \end{equation*}
\end{proof}

\begin{remark}
    Für eine normierte Algebra $\Aa$ ist die Multiplikation als Abbildung von $\Aa \times \Aa \to \Aa$ stetig, wobei $\Aa \times \Aa$ mit der Produkttopologie versehen ist.
\end{remark}

\begin{definition}
    Sei $\Aa$ eine Algebra.
    \begin{itemize}
        \item $e \in \Aa$ heißt Einselement\footnote{Wie man in der Algebra oft zeigt, bedingt die Existenz von Einselement bereits die Eindeutigkeit.}, falls $a \cdot e = a = e \cdot a$ für alle $a \in \Aa$.
    \end{itemize}
    Sei nun $e \in \Aa$ ein Einselement.
    \begin{itemize}
        \item Wir definieren
        $$ \Inv(\Aa) \coloneqq \{ a \in \Aa \mid \exists b \in \Aa : a \cdot b = e = b \cdot a \}. $$
        \item Sei $a \in \Aa$, so heißt
        $$ \rho_{\Aa}(a) \coloneqq \{ \lambda \in \C \mid (a - \lambda e) \in \Inv(\Aa) \} $$
        die \emph{Resolvente} von $a$. Weiters definieren wir durch $\sigma_{\Aa}(a) \coloneqq \C \setminus \rho_{\Aa}(a)$ das \emph{Spektrum} von $a$. Nach Kontext schreiben wir auch nur $\rho(a)$ und $\sigma(a)$.
        \item Ist $a \in \Aa$, so heißt $r_{\Aa}(a) \coloneqq \sup_{\lambda \in \sigma(a)} \vert \lambda \vert$ der \emph{Spektralradius} von $a$. Ist das Spektrum leer, so setzen wir das Supremum nach Konvention gleich $0$.
    \end{itemize}
\end{definition}

\begin{example}
    Für einen Hilbertraum $H$ betrachte $L_b(H)$ und sei ${}^*$ die Operation des Konjugierens eines Operators. Dann ist $(L_b(H), {}^*)$ sogar eine $C^*$-Algebra.

    Sei $P \in L_b(H)$ eine nichttriviale orthogonale Projektion, so gilt $P^2 - P = 0$. Nach dem Spektralabbildungssatz folgt $\sigma(P^2 - P) = \{ 0 \}$. Also ist $z \in \sigma(P)$ genau dann, wenn $z^2 - z = 0$, womit $\sigma(P) \subseteq \{ 0, 1 \}$ folgt. Da $P$ nichttrivial ist gibt es $x \in \ker P, y \in \ran P$ mit $x, y \neq 0$, was $P \notin \Inv(L_b(H))$ und $P - I \notin \Inv(L_b(H))$ impliziert und damit $\{ 0, 1 \} \subseteq \sigma(P)$.
\end{example}

\begin{lemma}
    Sei $\Aa$ eine Algebra mit Einselement und $a \in \Inv(\Aa)$, so gilt
    $$ \sigma(a^{-1}) = \{ 1/z \mid z \in \sigma(a) \}. $$
\end{lemma}

\begin{definition}
    Sei $\Aa$ eine normierte Algebra und $e \in \Aa$ ein Einselement. Gilt $\Vert e \Vert = 1$, so heißt $e$ \emph{normiert}. In dem Fall nennen wir $e$ auch \emph{Eins} und reden von einer Algebra mit Eins.
\end{definition}

\begin{remark}
    Sei $\Aa$ eine normierte Algebra mit Einselement, dann folgt, wegen $\Vert x^2 \Vert \leq \Vert x \Vert^2$ für alle $x \in \Aa$, stets $1 \leq \Vert e \Vert$.
\end{remark}

\begin{example}
    Ist $(\Aa, \Vert \cdot \Vert)$ eine Algebra mit Eins und definiert man für $t \geq 1$ eine Norm $\Vert \cdot \Vert_t : x \mapsto t \Vert x \Vert$, so ist $(\Aa, \Vert \cdot \Vert_t)$ eine normierte Algebra mit Einselement und $\Vert e \Vert_t = t$.
\end{example}

\begin{remark}
    Sei $\Aa$ eine Banachalgebra mit Eins. Dann gilt:
    \begin{itemize}
        \item $\Inv(\Aa)$ ist eine offene Teilmenge von $\Aa$.
        \item Für $a \in \Aa$ ist $\rho(a)$ als Teilmenge von $\C$ offen.
        \item Für $a \in \Aa$ ist $\sigma(a) \subseteq \C$ abgeschlossen. Weiters gilt sogar stets $\sigma(a) \subseteq K^\C_{\Vert a \Vert}(0)$, womit das Spektrum stets kompakt ist. Weiters ist das Spektrum stets nichtleer.
        \item Es gilt
        $$ r(a) = \lim_{n \to \infty} (\Vert a^n \Vert)^{1/n} = \inf_{n \in \N} (\Vert a^n \Vert)^{1/n}. $$
        \item Ist $\Aa$ sogar eine $C^*$-Algebra mit Eins, $x \in \Aa$ und gilt $x x^* = x^* x$, dann gilt sogar $r(x) = \Vert x \Vert$.
    \end{itemize}
\end{remark}

\begin{example}
    Betrachte $L_b(\C^2)$, wobei $\C^2$ als Hilbertraum aufzufassen ist und setze
    $$ J \coloneqq \begin{pmatrix}
        0 & 0 \\ 1 & 0
    \end{pmatrix}. $$
    Dann ist $\sigma(J) = \{ 0 \}$, womit $r(J) = 0 \neq 1 = \Vert J \Vert$.
\end{example}

\begin{remark} {\ }
    \begin{enumerate}
        \item Seien $\Aa, \mathcal{B}$ Algebren mit Einselement und $\varphi : \Aa \to \mathcal{B}$ ein Homomorphismus mit $\varphi(e_{\Aa}) = e_{\mathcal{B}}$. Dann folgt $\varphi(\Inv(\Aa)) \subseteq \Inv(\mathcal{B}).$ Für $a \in \Aa$ folgt $\rho_{\Aa}(a) \subseteq \rho_{\mathcal{B}}(\varphi(a))$, sowie $\sigma_{\Aa}(a) \supseteq \sigma_{\mathcal{B}}(\varphi(a))$.
        
        \item Sei $\Aa$ Algebra mit Einselement, $\mathcal{B}$ eine Algebra und $\varphi : \Aa \to \mathcal{B}$ ein surjektiver Homomorphismus. Dann hat auch $\mathcal{B}$ ein Einselement.
        
        \item Seien $\Aa, \mathcal{B}$ $C^*$-Algebren mit Einselement und $\varphi : \Aa \to \mathcal{B}$ ein ${}^*$-Homo\-morphismus mit $\varphi(e_{\Aa}) = e_{\mathcal{B}}$. Dann folgt $\Vert \varphi \Vert = 1$, da für $a \in \Aa$ gilt
        \begin{align*}
            \Vert \varphi(a) \Vert^2 &= \Vert \varphi(a^*) \varphi(a) \Vert = r(\varphi(a^* a)) = \sup_{\lambda \in \sigma(\varphi(a^* a))} \vert \lambda \vert \leq \sup_{\lambda \in \sigma(a^* a)} \vert \lambda \vert = r(a^* a) = \Vert a \Vert^2.
        \end{align*}

        \item Sei $\Aa$ eine Algebra mit Einselement $e$. Sei $\mathcal{B} \subseteq \Aa$ eine Unteralgebra mit $e \in \mathcal{B}$. Dann folgt $ \Inv(\mathcal{B}) \subseteq \Inv(\Aa) \cap \mathcal{B} $ und $\rho_{\mathcal{B}}(b) \subseteq \rho_{\Aa}(b) $ für alle $b \in \mathcal{B}$.
        
        \item Sei $\Aa$ eine Banachalgebra mit Eins $e$. Sei $\mathcal{B} \subseteq \Aa$ eine Banachunteralgebra mit $e \in \mathcal{B}$. Zunächst ist $\Inv(\mathcal{B})$ offen als Teilmenge von $\Inv(\Aa) \cap \mathcal{B}$, nach Definition der Spurtopologie. Sei nun $b \in \cl_{\Inv(\Aa) \cap \mathcal{B}}(\Inv(\mathcal{B})) = \cl_{\Aa}(\Inv(\mathcal{B})) \cap \Inv(\Aa) \cap \mathcal{B}$. Dann existiert $b^{-1} \in \Aa$. Gleichzeitig gilt $b = \lim_{n \to \infty} b_n$, wobei $b_n \in \Inv(\mathcal{B}) \subseteq \Inv(\Aa)$, womit sogar $b_n^{-1} \in \mathcal{B} \subseteq \Aa$. Da $\mathcal{B}$ abgeschlossen ist folgt sogar $b^{-1} \in \mathcal{B}$, und damit sogar $b \in \Inv(\mathcal{B})$. Also ist $\Inv(\mathcal{B})$ abgeschloffen als Teilmenge von $\Inv(\Aa) \cap \mathcal{B}$.
        
        Ebenso sieht man dass $\rho_{\mathcal{B}}(b) \subseteq \rho_{\Aa}(b) $ abgeschloffen ist.
    \end{enumerate}
\end{remark}

